%------------------------Thesis Master------------------------------------------------------------%
% Thomas Griffiths 2013-07-16 tmg994@uowmail.edu.au                                               %
% If you get stuck read my comments, here and in the preamble (thesispreamble.tex)                 %
% hopefully they can help you find where your answers will be.                                    %
% I highly reccomend making friends with Google, it knows quite a bit.                            %
%-------------------------------------------------------------------------------------------------%

\documentclass[12pt,oneside,english,]{book}

%------------------------ Preamble and bibliography resources
%--------------- UOW Thesis Preamble ----------------------------------------%
% Thomas Griffiths 2013-07-16 tmg994@uowmail.edu.au                          %
%                                                                            %
% I encourage you to read the documentation for each of the packages below,  %
% they contain instructions for implementation and examples of their use.    %
% If you get stuck read my comments, hopefully they can help you find where  %
% your answers will be. I highly reccomend making friends with señor Google, %
% he knows quite a bit. The wikibook on LaTeX is also very helpful:          %
% http://en.wikibooks.org/wiki/LaTeX/                                        %
% The packages I've loaded here are the bare basics. and mainly deal with    %
% formatting, captions and things. There are thousands of packages out there %
% for all the disciplines and formatting you might need. Google what you're  %
% looking for and the keywords 'LaTeX' and 'package', You'll probably find   %
% what you're looking for. I encourage you to look on ctan, www.ctan.org/‎,   %
% for packages that might be relevant to your degree. If you're in science I % 
% can recomend the siunitx and the chemmachro (or mhchem) package. They make %
% it really easy to typeset chemical equations, any quantity with units and  %
% scientific notation.                                                       %
%                                                                            %
%----------------------------------------------------------------------------%

\usepackage{geometry}
	\geometry{a4paper,inner=4cm, outer=2cm, top=3cm, bottom=2cm}
	% Dimensions from UOW thesis guidelines.
	\pdfpagewidth=\paperwidth 
	\pdfpageheight=\paperheight
	% This acts as a failsafe to ensure things aren't stretched or moved when it's finally printed as a PDF.

%\usepackage[parfill]{parskip} 
% Activate to begin paragraphs with an empty (return) line, comment out the indent below if you chose the return line option.

\setlength{\parindent}{4ex}	% Sets the length of the paragraph indent. Current setup has a an indent. Disable this if you activate the return line above.
	
\usepackage{setspace}
% Double or one and a half spacing.
	
\usepackage{graphicx}
	\DeclareGraphicsRule{.tif}{png}{.png}{`convert #1 `dirname #1`/`basename #1 .tif`.png}
% Graphics. Remove me and you won't have any figures, and that would be very boring.

\usepackage[usenames,dvipsnames,svgnames,table]{xcolor}
% Adds the ability to make coloured text and lines throughout the document. See documentation for xcolor.

%-------------------- Tables, figures and captions
\usepackage[font={small},labelfont={bf},margin=4ex]{caption}
% Makes bold labeled and smaller font captions. Must be loaded before the longtable package to avoid conflicts! 

\usepackage{longtable} 
% Long tables (more than one page). Different headers and footers for beginning and end pages, etc.

\usepackage{afterpage} 
% Make a longtable start on the next clear page, but fills the previous one with text first (no random gaps in the text-from long tables anymore! Man, the day I discovered this...)

\usepackage{booktabs} 
% Nice looking tables and lines in tables

\usepackage{multirow} 
% Entries in tables over multiple rows

\usepackage{lscape} 
% Pages in landscape

\usepackage{pdflscape} 
% Landscape pages also rotated in the pdf

\usepackage{wrapfig} 
% Allows figures that don't take up the entire width of the page, wrapping the text around the figure

%\usepackage[position=top,singlelinecheck=false,captionskip=4pt]{subfig} 
% Multiple figures in an individual figure. Fig. 1 a, b, c, etc. each with, or without, it's own individual caption, and with a global caption for all sub figures.

%-------------------- Special symbols and fonts
\usepackage{amssymb} 
% Maths symbols

%-------------------- Document sections, headers, footers, and bibliography
\usepackage{fancyhdr}											
% for creating different headers and footers

%-------------------- Bibliography
\usepackage[backend=biber,articletitle=true,style=chem-rsc,doi=false]{biblatex}
% This is the package that lets you create a bibliography. I recommend reading the biblatex documentation to understand all the options i've specified here. BibLaTeX was created to replace BibTeX. It has lots of extra fields and options. I'm also using the biber backend here rather than the default, it copes with unicode and so can deal with accented characters easily.

% Currently this is set up to use RSC style references with article titles displayed.

% Traditionally you would use BibTeX, a special build of TeX, the newer biblatex package is a more powerful bibliograpy management tool for LaTeX. You can make multiple chapter based bibliographies, footnote bibliographes, sort your references by date, author, order cited, essentially by any bit of citation data you happen to have. You can also have a seperate library with a differnet format for say books and articles. Or if you're a PhD student, the thesis references and your publications.

%\usepackage[numbers,super,comma,sort&compress]{natbib}				
	%\setcitestyle{square}										
	% places citations in square brackets to helps to distinguish between powers and citations
%This is the old natbib package that meshes with bibtex (rather than using the newer biblatex). It's here mainly for legacy purposes. Try to shift to biblatex if you can, it is cleaner in it's implementation and creating a custom citation style is easier then with bibtex.

\usepackage[unicode=true,colorlinks=true,linkcolor=black,citecolor=black,urlcolor=black,breaklinks=true]{hyperref}
% The hyperref package allows you to have clickable links in your pdf. It also allows you to have the mailto link associated with your name. It can be  a bit finicky, so load it last.

%-------------------- Command renewals, New commands etc.
\renewcommand{\thefootnote}{\alph{footnote}}							
%letters for footnotes instead of numbers to avoid confusion with references.

\addbibresource{Your_Bibliography.bib}

%------------------------Main Document--------------------------
\begin{document}
    \onehalfspace
	
%-------------- Title Page
% Title page info (see uowthesistitlepage package)
\title{A Pretty Swish Title} 
\author{Average J. Blow} % Full name, and any degrees held.
\date{January 1901} % Month Year, alternatively use the \today macro for Month dd, yyyy.
\degree{That Degree You're Studying} % Write it in full: e.g. Bachelor of Science Medicinal Chemistry Honours
\supervisor[2]{Dr. J Bloggs \& Professor J. Citizen}% The optional argument (default 1) in square brackets is the number of supervisors. In the Curly braces list your supervisor(s) seperated by commas.
\school{Your School} % e.g Chemistry

\maketitle
%-------------- Front Matter
\declaration
\pagestyle{fancy} \frontmatter
% These \phantomsection are to ensure that the hyperref package hyperlinks to the correct page in the electronic pdf. If you turn hyperref off they don't do anything so they can just stay here.
\phantomsection\addcontentsline{toc}{chapter}{Abstract}
\chapter*{Abstract} % Starred chapter=chapter with no number.
\chapter*{Acknowledgments}

\cleardoublepage
\phantomsection \pdfbookmark[0]{Contents}{Contents} % These \phantomsection are to ensure that the hyperref package hyperlinks to the correct page. If you turn hyperref off they don't do anything so they can just stay here
\tableofcontents

% A list of Figures, uncomment to activate.
%\cleardoublepage
%\phantomsection \label{listoffigures}\addcontentsline{toc}{chapter}{List of Figures}
%\listoffigures

% A list of Schemes, uncomment to activate.
%\cleardoublepage
%\phantomsection \label{lisrofschemes}\addcontentsline{toc}{chapter}{List of Schemes} 
%\listofschemes

% A list of Tables, uncomment to activate.						
%\cleardoublepage
%\phantomsection \label{listoftables}\addcontentsline{toc}{chapter}{List of Tables} 
%\listoftables

%-------------- Chapters
\mainmatter
\pagestyle{fancy}
\chapter{Introduction}
\section{The Title of Your First Section}
\subsection{A Subsection---for Clarity}
Lorem ipsum dolor sit amet, consectetur adipiscing elit.\cite{pericles} Integer ut congue lectus. Nullam dapibus scelerisque diam, ac convallis dolor convallis non. Sed nec lectus nec sapien interdum commodo nec quis elit. Sed neque augue, pulvinar id imperdiet id, ullamcorper sed sem. Duis pulvinar blandit erat, quis suscipit erat venenatis sit amet. Aenean varius aliquam dignissim. Sed tempus consequat sapien et bibendum. Class aptent taciti sociosqu ad litora torquent per conubia nostra, per inceptos himenaeos. Maecenas arcu ligula, tincidunt non volutpat nec, luctus ac justo. Donec euismod egestas leo, vel convallis mi accumsan sed. Aliquam erat volutpat. Maecenas volutpat lacinia justo, dictum iaculis enim scelerisque in. Vestibulum consequat augue in nisl luctus eget ultrices sem blandit. Donec mi risus, rutrum at tempus in, dapibus vel est. Sed eu ullamcorper velit.\cite{pericles,thekingsenglish}

Duis malesuada ultrices rutrum. In posuere sem dapibus urna accumsan id sollicitudin turpis iaculis. Sed fermentum metus vel dui vulputate consectetur. Nunc nec ante nisi, ac ultricies metus. Nullam sed dui vitae metus vehicula adipiscing. Nunc facilisis tortor neque. Praesent ac nulla at odio lacinia egestas. Nam facilisis vehicula pretium.

Fusce tempus libero vitae leo cursus cursus. Suspendisse iaculis dignissim placerat. Vivamus ornare, lectus in placerat cursus, quam ligula congue risus, a facilisis leo est ut massa. Nulla in metus id eros rutrum aliquam. Cras id velit ut enim iaculis adipiscing facilisis sit amet lorem. 1234567.90 Ut id neque id turpis facilisis tempus sit amet nec urna.\cite{bioluminescence:ch1} Ut volutpat magna et dolor elementum eget elementum lorem eleifend. Pellentesque mollis convallis est a tincidunt. Nunc egestas, ante a blandit iaculis, risus orci auctor risus, in pellentesque est sapien id mi. Nunc faucibus porttitor tincidunt.\cite{atkins_inorgchem,atkins_physchem} Fusce eu faucibus dolor. Mauris quis arcu metus, at malesuada lacus. Nam lectus lacus, sodales nec tempus eu, pretium id risus. Cras dignissim aliquet laoreet. Morbi sed arcu id tellus tempor consectetur.\cite{bioluminescence:apA,hori1973}

\subsection{Another Subsection}
Aenean convallis, ante quis convallis viverra, urna lorem venenatis odio, ut bibendum nisi mauris a mi. Nam laoreet arcu a eros dignissim elementum. Class aptent taciti sociosqu ad litora torquent per conubia nostra, per inceptos himenaeos. Proin vitae magna ligula, id consequat felis. Integer rutrum ante vitae tellus convallis volutpat in vitae odio. Vivamus eu turpis et turpis euismod faucibus. Suspendisse at turpis sit amet magna auctor fringilla. Vestibulum quis nulla sed dui pellentesque dapibus. Cum sociis natoque penatibus et magnis dis parturient montes, nascetur ridiculus mus. Maecenas id odio nisl. Etiam facilisis elit in urna hendrerit ac iaculis mauris lacinia. Fusce congue ultricies nulla vel elementum. Phasellus suscipit vestibulum tortor at semper. Etiam eu sagittis augue.\cite{gaussian09,bioluminescence,clayden_orgchem}
\begin{figure}
\begin{center}
\includegraphics[height=4cm]{figures/example.jpg}
\caption[Chemistry cat.]{This is chemistry cat. Here he serves to demonstrate a figure in a LaTeX document, complete with caption. Notice his suave bow tie, but also how that he has forgotten to label his solutions.}
\end{center}
\end{figure}
Duis malesuada ultrices rutrum. In posuere sem dapibus urna accumsan id sollicitudin turpis iaculis. Sed fermentum metus vel dui vulputate consectetur. Nunc nec ante nisi, ac ultricies metus. Nullam sed dui vitae metus vehicula adipiscing. Nunc facilisis tortor neque. Praesent ac nulla at odio lacinia egestas. Nam facilisis vehicula pretium.

Praesent sed diam arcu, quis aliquam libero. Pellentesque a nulla vulputate lectus rhoncus vulputate ut in urna. Quisque sapien sem, convallis ac porta sit amet, aliquet sit amet eros. Etiam quis lorem ligula. Mauris elit sapien, ultrices nec vulputate ut, porta eu leo. Ut consequat accumsan commodo. Sed ut ultricies ante.\cite{mithani2012}

\chapter{Your Second Chapter}

\chapter{Your Third Chapter}

\chapter{Your Fourth Chapter}

\chapter{Your Fifth Chapter}


%-------------- Bibliography
\cleardoublepage
\phantomsection	\addcontentsline{toc}{chapter}{Bibliography}										
\printbibliography

%-------------- Apendicies
\cleardoublepage
\noappendicestocpagenum % Removes the number in the toc next to the 'appendices' entry. Each individual appendix still has a page number in the toc.
\begin{appendices}
\chapter{Appendix One}
\section{The title of the first section}
The appendices work exactly the same way as chapters, they are numbered with letters rather than numbers though.
\chapter{Appendix Two}
\section{The title of the first section}
\end{appendices}

\end{document}  

%---------------------------Notes-------------------------------

%-------------- 
% Floats and centreing for tables and figures \begin{table,scheme,figure} go in the body text, *then* include the file without any floats or centreing, closing with \end{table,scheme,figure} to ensure correct float placement. This is only if you are calling in your tables as seperate files, not leaving them in the text. Also note the american spelling of center (not the standard AU centre).

% e.g.

% Figure or Scheme (caption after)

% \begin{figure}
% \begin{center}
% \input{figures/name_of_figure.tex}
% \caption[Short caption for table of figures]{Long caption for text body}
% \end{center}
% \end{figure}

% Table (caption before)

% \begin{table}
% \begin{center}
% \caption[Short caption for table of tables]{Long caption for text body}
% \input{figures/name_of_figure.tex}
% \end{center}
% \end{table}

%-------------- 
% The difference between \input and \include is thus: \include creates an aux file for the included file, \input does not. \include is good for large files, like a chapter, since LaTeX won't reread the file if no changes have been made. \input is better for small inputs like data tables, since it doesn't create an aux file for the inputted file, it is read each and every time LaTeX is run. Additionally you *CAN* nest \input commands. you *CANNOT* nest \include commands.

%-------------- 








